\chapter{Metodologia de avaliação}
\label{cap:metodologia}

% wiki: LLM Stage Design (Testing and evaluation), concepts/Precision Metrics,
% benchmarks/Racebench; project-src/docs/evaluation-protocol.md

Este capítulo define como a ferramenta será avaliada: as perguntas de pesquisa,
o gabarito contra o qual se mede, as métricas, o protocolo experimental e as
ameaças à validade.

\section{Perguntas de pesquisa}
\label{sec:qp}

\begin{description}
  \item[QP1] A etapa estática encontra todos os defeitos anotados no Racebench,
    isto é, atinge recall de 100\%?
  \item[QP2] A triagem por \llm\ reduz o esforço de revisão sem que nenhum
    defeito real seja classificado como inviável?
  \item[QP3] Qual a contribuição de cada técnica de \textit{prompt}? Essa pergunta
    é respondida por uma ablação no formato de \citeonline{llift}.
  \item[QP4] O \llm\ propõe correções corretas para os defeitos encontrados?
\end{description}

\section{Benchmarks e gabarito}
\label{sec:benchmarks}

\subsection{Unidade de contagem}

Toda taxa de recall divide por um número de defeitos, e a literatura usa números
diferentes para os mesmos programas sem dizer por quê
(Seção~\ref{sec:rel-comparacao}). Este trabalho adota explicitamente
uma tripla anotada como um defeito. Medidas sobre os 31 casos simples
do Racebench, as unidades possíveis dão resultados bem diferentes
(Tabela~\ref{tab:unidades}).

\begin{table}[htb]
  \centering
  \caption{Contagem do gabarito do Racebench conforme a unidade}
  \label{tab:unidades}
  \begin{tabular}{@{}lcc@{}}
    \toprule
    Unidade & Defeitos & Armadilhas \\
    \midrule
    por instância de tripla & 48 & 38 \\
    por par (caso, variável) & 33 & 29 \\
    por caso & 31 & 31 \\
    \bottomrule
  \end{tabular}
  \fonte{Elaborada pelos autores a partir das anotações de \citeonline{racebench}.}
\end{table}

A unidade por instância é a mais fina, e portanto a mais difícil de inflar; é a
usada pelo BMC4AV, o que preserva a comparação; e as unidades mais grossas
escondem estrutura: um mesmo caso chega a ter quatro defeitos distintos
numa só variável.

\subsection{Leitura e correção do gabarito}

As anotações do Racebench usam ao menos quatro gramáticas incompatíveis e contêm
erros: um leitor estrito encontra apenas 28 defeitos, sem acusar erro. Por isso
o leitor de gabarito precisa tolerar as variações, e sua contagem é validada
contra uma contagem manual dos casos mais difíceis antes de qualquer medição.

Doze anotações são corrigidas antes da comparação, por uma errata em que cada
correção traz sua evidência: seis números de linha (cinco deles num único caso cujas anotações apontam para linhas sem nenhum acesso) e seis tipos de acesso trocados.
Nenhuma correção altera as contagens. Sem elas, uma comparação exata com o texto
distribuído daria recall zero para um defeito que a ferramenta encontra na linha
certa.

\subsection{Regra de casamento}

Um candidato reportado corresponde a uma tripla anotada quando está no
mesmo programa, trata da mesma posição de memória e tem as mesmas linhas
ordenadas $(A_1, B, A_2)$ que a anotação corrigida. Três escolhas são
deliberadas:

\begin{itemize}
  \item o tipo de acesso não precisa coincidir, porque as linhas já distinguem
    todas as 86 anotações sem colisão e 20 acessos anotados estão em linhas que
    leem e escrevem a variável;
  \item o nome da variável é comparado depois de resolver apelidos, pois 12
    anotações nomeiam a posição por um ponteiro, um elemento de vetor ou um campo de estrutura;
  \item o recall é medido apenas sobre triplas; os pares recebem uma medida
    própria de cobertura, publicada ao lado.
\end{itemize}

\subsection{Conjunto de triagem}

A etapa com \llm\ é desenvolvida e medida, até a integração com a etapa
estática, sobre 20 registros de contexto montados à mão a partir do
Racebench: 10 defeitos anotados e 10 armadilhas, organizados de modo a formar
seis pares adversariais, isto é, um defeito e uma armadilha que diferem em um
único detalhe, como o estado de mascaramento. Os pares existem para que o modelo não
acerte pela aparência do código. Um caso cujas anotações descrevem uma versão
diferente do programa distribuído foi excluído do conjunto.

\subsection{Suíte de programas reais}

A suíte de 18 programas (Seção~\ref{sec:rel-benchmarks}) é usada na etapa final.
Seu gabarito registra, por variável, 45 violações verdadeiras e 14 falsas, o
que também a torna um conjunto de triagem.

\section{Métricas}
\label{sec:metricas}

\begin{description}
  \item[Portão de recall] Todos os defeitos anotados precisam chegar ao relatório
    e nenhum pode ser classificado como provavelmente inviável. Uma falha aqui
    é um defeito de projeto, não um ajuste.
  \item[Rejeição de armadilhas] Fração das armadilhas classificadas nos grupos
    baixos. É a medida de precisão da triagem e, ao contrário da precisão, não
    melhora descartando candidatos.
  \item[\textit{Inspection Ratio}] Fração da lista ordenada que um revisor
    precisa ler até encontrar todos os defeitos \cite{reducing-fp}. Empates de
    grupo e confiança são resolvidos pelo pior caso, colocando o defeito real por
    último.
  \item[Consistência] Concordância entre execuções repetidas do mesmo candidato.
    \citeonline{llift} mostram que a decomposição e a autovalidação a melhoram
    independentemente da acurácia.
  \item[Taxas de reparo] Taxa sintática (um analisador sintático aceita a correção) e taxa
    lógica (a correção equivale a uma escrita à mão), como em
    \citeonline{skipanalyzer}, seguidas da re-verificação.
\end{description}

No Racebench, apenas a viabilidade entra nessas métricas. O gabarito anota
fatos do programa (a intercalação existe), e muitos de seus defeitos leem
a variável compartilhada para uma variável local que nunca é usada. Um modelo
que classifica esses casos como benignos está raciocinando corretamente sobre o
código, e pontuar essa resposta contra o gabarito puniria o acerto. A pergunta
de nocividade continua sendo feita e registrada, e é avaliada na suíte de
programas reais, em que o código tem efeito observável.

\section{Protocolo experimental}
\label{sec:protocolo}

\subsection{Ablação}

A ablação acrescenta um componente de \textit{prompt} por linha, sempre sobre os
mesmos candidatos: \textit{prompt} simples; mais regras do domínio; mais pedidos
progressivos de contexto; mais decomposição da tarefa; mais autovalidação. Como
os pedidos de contexto só trazem informação nova quando respondidos a partir do
código real, a linha correspondente é medida depois da integração com o
resolvedor da etapa estática.

Na ablação do LLift, o \textit{prompt} simples já perdia a maioria dos defeitos,
e o recall era a variável que se movia. Se, neste domínio, o recall já partir de
100\%, as variáveis dependentes da ablação passam a ser a rejeição de armadilhas
e o \textit{Inspection Ratio}.

\subsection{Registro e reprodutibilidade}

Todo número é registrado com o modelo, a data e a configuração de
\textit{prompt} que o produziram, pois as comparações de modelos da literatura
envelhecem rapidamente. As respostas do modelo são guardadas em cache,
indexadas pelo candidato, pela configuração e pelo modelo, de modo que uma
mesma medição possa ser recalculada sem nova chamada. O custo em
\textit{tokens} e em tempo é medido por candidato. Uma linha da ablação só é
considerada medida se cobrir todos os candidatos previstos.

\section{Ameaças à validade}
\label{sec:ameacas}

\begin{itemize}
  \item Amostra pequena. Vinte candidatos e uma amostra por candidato
    tornam diferenças de um ou dois candidatos indistinguíveis de ruído; a
    medição de consistência com várias amostras é o que permite afirmar uma
    diferença.
  \item Ajuste sobre o conjunto de teste. Quando o \textit{prompt} é
    corrigido depois de observar erros nos mesmos 20 casos, o número seguinte é
    otimista. Por isso as versões anteriores e posteriores são reportadas lado a
    lado.
  \item Gabarito com erros conhecidos. A errata corrige o que pode ser
    verificado no código; o rótulo de um caso cujas anotações descrevem outra
    versão do programa não pode ser corrigido sem decidir de novo o gabarito.
  \item População diferente após a integração. Com a etapa estática
    completa, o \llm\ recebe apenas o que o solver não decidiu, que são os casos
    mais difíceis; os números obtidos sobre os 20 registros montados à mão não
    se transferem automaticamente.
  \item Benchmarks de arquivo único. Todos os programas avaliados são um
    único arquivo; a entrada com vários arquivos é uma capacidade da ferramenta
    exercida fora do gabarito.
  \item Premissas não conservadoras. A exclusão de reentrância
    (Seção~\ref{sec:premissas}) é uma perda potencial de recall registrada, não
    medida.
\end{itemize}
